% This file should be rename with the speaker's last name. (e.g : smith)

% Write the author names in the order you wish them to appear
% and underline the one who will give the talk.

% Only the speaker's email address is required.

\begin{abstract_online}{Gaussian Elimination with Parameters}{%
        \underline{Aharon Naiman}$^{1}$}{[naiman@jct.ac.il]
        }{%
        $^1$ Department of Applied Mathematics, Jerusalem College of
    Technology, Jerusalem, Israel}
    
Basic mathematics courses, at all levels, involve \textit{many}
opportunities to include CAS packages.
%like \emph{Mathematica} [2], Maple [1], REDUCE [3], Sage [4],
%amongst many others.  Computer algebra packages
Such systems assist with the preparation of:
%
\begin{itemize}
    \item classroom slides/notes,
    \item individualized homework assignments,
    \item in-class, randomized quizzes,
    \item class projects,
    \item extra-credit, further reading,
    \item final examinations,
    \item etc.
\end{itemize}
%
In this talk we discuss an aspect which affects all of the areas
above, i.e., that of solving Gaussian elimination with parameters, in
particular for the teaching of basic, first-semester linear algebra.

Right from the beginning of the semester, students are shown how to
perform row reduction.  As we know, they need to show that there are
either no solutions, one unique solution (and what it is), or an
infinite number of solutions (and what they are).  Are the standard
functions of the available packages prepared to show these three
possibilities?

The linear systems are then ``complicated'' by including input
parameters.  The students need to continue to solve these systems, and
specify, based on the input parameters, which of the three
possibilites above pertains.  Again, do the standard, available
functions supply all of the necessary solutions?  As we shall show,
not all solutions and special cases are covered.

Three approaches are presented using \emph{Mathematica} [2], including
one which gets back to basic, row reduction.  This last one is
particularly useful, as it does all of the calculations itself (\`{a}
la Computer-Based Maths [1]), step-by-step, so the student misses
nothing, and is nonetheless not bogged down with a myriad of
arithmetic calculations.  This leaves more time for
\emph{understanding} the matrix (or individual vectors), as well as
applications of solutions of linear systems.

We compare the approaches, demonstrating that some have more
satisfying results than others, handling all special cases (and
\emph{not} unnecessary ones).  We show that one of the approaches
delays the need for handling special cases of parameter values along
the way, obviating the need for students to recall these special cases
until the end.

We end off with applying a final approach to most of the exercises
posed in the remainder of the linear algebra course.

\textbf{Keywords} \newline{} linear algebra; education; automated
Gaussian elimination with parameters; \emph{Mathematica}

\textbf{References} \newline{}
\emph{Computer-Based Maths} at \texttt{computerbasedmath.org} \newline{}
%\emph{Maple} at \texttt{www.maplesoft.com/products/Maple} \newline{}
\emph{Mathematica} at \texttt{www.wolfram.com/mathematica} \newline{}
%\emph{REDUCE} at \texttt{reduce-algebra.com} \newline{}
%\emph{Sage} at \texttt{www.sagemath.org} \newline{}
\end{abstract_online}
    
